\section{Cómo usar el modelo aprendido}

\subsection{Método 1: planificación basada en el modelo (Planning)}

Con un modelo de dinámica disponible, se puede buscar directamente en el modelo la secuencia de acciones óptima:

$$\max_{\mathbf{a}_{t:t+H}} \sum_t r(\mathbf{s}_t, \mathbf{a}_t)$$

\begin{figure}[H]
\centering
\includegraphics[width=0.9\textwidth]{slides-images/slide-15.jpg}
\caption{Método 1b: planificación mediante optimización sin gradientes por muestreo\protect\footnotemark}
\end{figure}
\footnotetext{Slides, página 15.}

Flujo concreto del algoritmo:
\begin{enumerate}
    \item Ejecutar alguna política (por ejemplo, una política aleatoria) para recolectar datos $\mathcal{D} = \{(\mathbf{s}, \mathbf{a}, \mathbf{s}')_i\}$
    \item Aprender el modelo $f_\phi(\mathbf{s}, \mathbf{a})$
    \item Muestrear iterativamente secuencias de acciones y elegir la mejor acción mediante la simulación hacia adelante con el modelo
\end{enumerate}

\subsubsection{Método de disparo aleatorio (Random Shooting)}

El método más sencillo consiste en muestrear aleatoriamente $N$ secuencias de acciones, hacer el rollout de cada secuencia en el modelo y elegir la primera acción de la secuencia con la mayor recompensa acumulada.

\subsubsection{CEM (Cross-Entropy Method)}

CEM es un método de optimización sin gradientes más avanzado:
\begin{enumerate}
    \item Inicializar una distribución de acciones (por ejemplo, una distribución gaussiana)
    \item Muestrear $N$ secuencias de acciones de la distribución
    \item Evaluar en el modelo la recompensa acumulada de cada secuencia
    \item Seleccionar las top-$K$ secuencias «de élite»
    \item Actualizar la media y la varianza de la distribución de acciones con las secuencias de élite
    \item Repetir durante varias iteraciones
\end{enumerate}

\begin{knowledgebox}{MPC (Model Predictive Control)}
En la ejecución real se suele adoptar el paradigma MPC: en cada paso solo se ejecuta la primera acción obtenida por la planificación y luego se vuelve a planificar desde el nuevo estado. Esto mitiga en parte el error acumulado del modelo en horizontes temporales largos.
\end{knowledgebox}

\subsubsection{Método 1a: planificación basada en gradientes}

Si el modelo es diferenciable, se puede optimizar calculando directamente el gradiente respecto a la secuencia de acciones. Sin embargo, en la práctica este método tiende a quedar atrapado en óptimos locales y exige que el modelo sea suficientemente suave.

\subsection{Método 2: usar el modelo para generar datos sintéticos}

\begin{figure}[H]
\centering
\includegraphics[width=0.9\textwidth]{slides-images/slide-25.jpg}
\caption{Dos vías de optimización de políticas Model-Based\protect\footnotemark}
\end{figure}
\footnotetext{Slides, página 25.}

Principales limitaciones de los métodos de planificación:
\begin{itemize}
    \item Alto costo de cómputo en tiempo de prueba (hay que optimizar en cada paso)
    \item Están limitados por el problema del horizonte corto
\end{itemize}

Dos enfoques para resolver el problema del horizonte largo:
\begin{enumerate}
    \item \textbf{Planificación con función de valor terminal}: planificación de horizonte corto + una función de valor aprendida como estimación del valor terminal
    \item \textbf{Métodos estilo Dyna}: usar el modelo para generar datos sintéticos que refuercen el entrenamiento de RL model-free
\end{enumerate}

\begin{importantbox}{Algoritmo Dyna}
Dyna es un paradigma clásico de RL model-based:
\begin{enumerate}
    \item Recolectar datos reales y entrenar o actualizar el modelo de dinámica
    \item Usar el modelo para generar datos sintéticos (synthetic rollouts)
    \item Mezclar los datos sintéticos con los reales y actualizar la política con un algoritmo de RL model-free
\end{enumerate}
Su ventaja central es aprovechar la capacidad de generalización del modelo para «amplificar» una cantidad limitada de datos reales en una gran cantidad de datos de entrenamiento.
\end{importantbox}

\subsection{Manejo del error del modelo}

\begin{warningbox}{Efecto acumulativo del error del modelo}
Los pequeños errores del modelo en cada paso se acumulan a lo largo del eje temporal. Si el horizonte de planificación es $H$ y el error por paso es $\epsilon$, el error total después de $H$ pasos puede llegar a $O(H \cdot \epsilon)$ o ser aún mayor (si la dinámica es inestable). Por ello, confiar ciegamente en el modelo lleva a que la política «explote» los defectos del modelo (model exploitation).
\end{warningbox}

Estrategias de mitigación:
\begin{itemize}
    \item \textbf{Ensamble de modelos (Ensemble)}: entrenar varios modelos y usar las predicciones de los distintos modelos para estimar la incertidumbre
    \item \textbf{Rollout de horizonte corto}: limitar el número de pasos del rollout en el modelo
    \item \textbf{Actualización continua del modelo}: actualizar el modelo constantemente a medida que se recolectan datos nuevos
\end{itemize}

\subsection{Resumen de la sección}

El modelo puede usarse de dos maneras: (1) planificación directa (optimizar la secuencia de acciones) y (2) generación de datos sintéticos para reforzar el RL model-free. Ambas requieren manejar adecuadamente el error del modelo.
